\section{Introduction}

Semantic-ID (\sid) tokenization has become a practical artifact boundary for
generative recommendation. Systems such as TIGER map items into discrete code
sequences so that a sequence model can generate candidate identifiers rather
than score every catalog item directly~\cite{rajput2023tiger}. Once exported,
that item-to-code mapping becomes the address space exposed to the generator.
Recent work has expanded this space quickly: residual quantization frameworks,
recommender-native encoding, collaborative tokenization, differentiable
tokenization, non-uniform capacity, qualified collisions, variable-length
interfaces, collision-aware evaluation, and drift-aware refreshes all change
the artifact that downstream users must inspect before generator
training~\cite{ju2025grid,liang2026resid,zhu2024cost,fu2026diger,wei2026card,hu2026quasid,zhang2026sidreliability,cheng2026capsid,feng2026dact}.
This paper treats those exports as a shared artifact interface: the goal is not
to reimplement the literature, but to make item-to-code mappings
inspectable under the same contract.

Reusable tooling has not caught up with this artifact boundary. A new
tokenizer repository may expose checkpoints, item mappings, codebooks, or only
intermediate feature files. Downstream users then need to answer basic
engineering questions before training a generator: do all catalog items have
valid codes, are many items aliased to the same full code, do prefixes reflect
user co-occurrence rather than only metadata taxonomy, are tail items allocated
enough address capacity, and does the prefix structure create large fan-out?
Today these checks are usually reimplemented ad hoc inside each paper. The
resulting evidence is difficult to compare because each method reports a
different mixture of codebook usage, downstream ranking, and qualitative
examples.

Aggregate ranking metrics remain necessary, but they are the wrong interface
for inspecting the exported address space itself. Prior recommender tooling has
argued for behavioral tests beyond NDCG-style aggregates~\cite{chia2022reclist};
\sid tokenizers add a more concrete resource problem: the item mapping should
be validated and profiled before a full generator consumes GPU time. Underused
codebooks, repeated full codes, behaviorally weak prefixes, tail compression,
and large prefix fan-out can be hidden by a downstream score, yet they are
direct properties of the reusable tokenizer artifact.

We present \tool, a diagnostic/interface resource for inspecting item-to-\sid{}
tokenizer artifacts before or alongside downstream recommendation evaluation.
\tool is deliberately neither a new tokenizer nor a RecBole-style coverage
benchmark. It asks a narrower and reusable question: given any tokenizer
artifact that can emit item-level code sequences, what can be said about its
capacity use, aliasing, behavioral alignment, head-tail allocation, and
structural cost before retraining a full generator?

The resource contribution is threefold. First, \tool defines a small adapter
contract for \sid assignments, item metadata, interaction histories, and
optional generator outputs, plus validators that reject incomplete or
ambiguous artifacts before metrics are reported. Second, it implements
D1--D5 mapping-level probes for utilization, aliasing, neighborhood alignment,
popularity allocation, and structural cost, with D6 churn support for
continual-tokenizer settings. Third, it provides worked examples across four
tokenizer artifact lines: GRID/RQ-KMeans-style and ReSID/GAOQ form the
same-item diagnostic contrast, while LETTER and LC-Rec test released
item-index adapters. References, controls, and mechanism probes calibrate the
diagnostics. The central finding is narrow but useful: learned addressability
and behavioral prefix alignment diverge, so they should be inspected
as separate artifact properties before method leaderboard evaluation~\cite{cikm2026resource}.
